\section{SF-E839. Five leading brackets already obstruct a common point
defect}\label{SF-E839:sf-e839.-five-leading-brackets-already-obstruct-a-common-point-defect}

Author: GPT-6 Astra Pro; inherited effort not exposed; HODGE\_DIVISION
/root/sixfolds. Date: 2026-11-21.

\textbf{Status:} complete new refinement checked in full by the author, /root
and /root/g34\_eightfolds. The root also read the entire canonical
conversion and both exact certificate reports. E838 is retained unchanged. The new necessary
condition is that every leading bracket pencil have rank at most four.
An exact nonsquare rank-five example, with identically zero determinant
pencil, proves strict new scope relative to E838 as well as E836 and
E837.

\subsection{1. The stronger filtered-matrix
lemma}\label{SF-E839:the-stronger-filtered-matrix-lemma}

Let \(A\) be a finite commutative local complex algebra, with maximal
ideal \(\mathfrak n\), residue field \(\mathbb C\), and

\[
 V=\mathfrak n/\mathfrak n^2,\qquad W=\mathfrak n^2/\mathfrak n^3,
 \qquad e=\dim V.
 \tag{E839.1}
\]

Suppose that a square matrix \(C\in\operatorname{Mat}_m(A)\) has
cokernel annihilated by elements \(p_1,\ldots,p_s\in\mathfrak n^2\)
whose classes in \(W\) have span of dimension \(\rho>e\). Then \(C\) is
invertible.

Here is the complete dimension argument, so the stronger hypothesis is
not hidden in a reference to E838. Split all unit pivots by invertible
row and column operations. If a nonempty residual square block \(C'\) of
size \(k\) remains, all its coefficients lie in \(\mathfrak n\). Choose
complex vector-space splittings modulo \(\mathfrak n^3\), and write
\(C'=C_1+C_2\), with coefficients in the respective lifts of \(V,W\).
Let

\[
 L:\mathbb C^k\longrightarrow V\otimes\mathbb C^k,\qquad
 L(v)=C_1v,\qquad a=\operatorname{rank}L\le k.
 \tag{E839.2}
\]

A constant input has no first-order image exactly when it belongs to
\(\ker L\). Linear inputs contribute through \(C_1\), and quadratic
inputs contribute nothing modulo \(\mathfrak n^3\). Therefore the pure
second-order part of the image is exactly

\[
 \operatorname{im}(C'\bmod\mathfrak n^3)\cap(W\otimes\mathbb C^k)
   =C_2(\ker L)+V\cdot\operatorname{im}L.
 \tag{E839.3}
\]

The hypothesis forces \(e\ge1\): if \(e=0\), nilpotence gives
\(\mathfrak n=\mathfrak n^2=0\) and hence \(W=0\), contrary to
\(\rho>e\). Its dimension is at most \((k-a)+ea\le ek\). Each of the \(\rho\)
independent leading annihilators times every constant target vector lies
in that image, providing a subspace of dimension \(\rho k>ek\). This is
a contradiction. Thus no residual block remains. Only the actual linear
span of these leading annihilators was needed; their ideal need not be
all of \(\mathfrak n^2\).

The strict inequality cannot simply be replaced by equality in this
abstract lemma. In the four-variable square algebra, multiplication by
\(l=t_1+t_2+t_3+t_4\) is not invertible, and its cokernel is annihilated
by the four quadrics \(lt_i\). These are independent: a relation has
coefficients \(a_i+a_j=0\) at every \(t_it_j\); using three distinct
indices gives \(a_i=0\) in characteristic zero. This is only a boundary
for the matrix lemma, not a simultaneous module-deformation example.

\subsection{2. The rank-at-most-four necessary
condition}\label{SF-E839:the-rank-at-most-four-necessary-condition}

Keep the formal geometric setting of E838. Namely
\(R=\mathbb C[[z_1,z_2,z_3,z_4]]\), the two constant rank-two bivectors
\(\pi_+,\pi_-\) have complementary supports, and the leading quadrics
\(F_1,\ldots,F_4\) of the regular generators
\(f_i=F_i+O(\mathfrak m^3)\) form an Artinian homogeneous regular
sequence. Put

\[
 Q=\operatorname{span}(F_1,F_2,F_3,F_4),\qquad
 A=R/(f_1,f_2,f_3,f_4),\qquad W=S_2/Q.
 \tag{E839.4}
\]

The complete filtered initial-form proof of E838 gives
\(\operatorname{gr}A=S/(F)\), \(e=4\), \(\dim W=6\), and length sixteen.
It also proves that the two leading bracket tuples
\((\{F_i,F_j\}_{\pi_\pm}\bmod Q)_{i<j}\) are independent, using the
nonzero Jacobian socle in this homogeneous complete intersection.

Consider an actual common point modification

\[
 0\longrightarrow A^n\longrightarrow M\longrightarrow R^r
       \longrightarrow A^m\longrightarrow0,\qquad m+n>0.
 \tag{E839.5}
\]

If \(M\) has flat first-order deformations for both multiplication
classes, E819 gives \(n=m>0\), and E836 gives the actual common
companion matrices

\[
 C_\pm B_{ij}=B_{ij}C_\pm=P^\pm_{ij}I_m,\qquad
 P^\pm_{ij}=\{f_i,f_j\}_{\pi_\pm}\bmod I.
 \tag{E839.6}
\]

For any complex \(a,b\), form \(C_0=aC_++bC_-\). Its cokernel is
annihilated by \(P^0_{ij}=aP^+_{ij}+bP^-_{ij}\). If their classes in
\(W\) span a space of dimension at least five, the lemma makes \(C_0\)
invertible. Then

\[
 (C_\pm C_0^{-1})P^0_{ij}=P^\pm_{ij}I_m.
 \tag{E839.7}
\]

Taking the second filtration quotient and applying a linear functional
nonzero on one nonzero \(\overline P^0_{ij}\) forces the constant part
of each \(C_\pm C_0^{-1}\) to be scalar. All six equations therefore
make both leading tuples scalar multiples of the same one, contrary to
their independence. This proves:

\textbf{Rank-five obstruction.} For every such common point modification
admitting both deformations, the linear map

\[
 T_{a,b}:\Lambda^2Q\longrightarrow S_2/Q,\qquad
 F_i\wedge F_j\longmapsto
       \{F_i,F_j\}_{a\pi_++b\pi_-}\bmod Q
 \tag{E839.8}
\]

has rank at most four for every \(a,b\in\mathbb C\). Hence a single
rank-five or rank-six pencil value excludes weak balanced semiregularity
in the actual coherent-sheaf setting of E819, with its specified
coherent finite translation/tensor action and rational algebraic \(B\)
shear.

There is no extra deformation for a pencil parameter and no base change
of an Ext group or module deformation to a square quotient. Only the two
necessary identities for the original module were combined.

\subsection{3. The exact residual rank
equations}\label{SF-E839:the-exact-residual-rank-equations}

On a standard chart of \(\operatorname{Gr}(4,S_2)\), choose the ordered
generator matrix with a fixed four-by-four identity pivot block. After
reducing the six bracket columns by those four generators, \(T_{a,b}\)
is an explicit six-by-six matrix with entries linear in \(a,b\). The
rank condition is equivalent to vanishing of all its thirty-six
five-by-five minors, identically in the pencil variables. Each such
minor has degree five in \(a,b\), so its six coefficients give at most

\[
 36\cdot6=216
 \tag{E839.9}
\]

regular polynomial equations on each chart.

The chart polynomials have degree at most fourteen in the quadratic
coefficients. Indeed each five-by-five quotient minor is a nine-by-nine
augmented minor: retain all four pivot rows, the chosen five other rows,
all four generator columns and the chosen five bracket columns. Its
determinant is unchanged by eliminating the pivot columns. The four
generator columns have coefficient degree one, and the five bracket
columns have coefficient degree two, giving total degree fourteen before
normalization; substituting the identity pivot block does not increase
that degree.

These rank conditions are independent of the chosen generator basis and
quotient basis because rank is invariant under invertible changes of
both. They therefore define an intrinsic closed residual locus within
the regular-quadratic locus. No independence of the 216 equations,
codimension count, sufficiency, or existence of a module on the residual
locus is asserted.

\subsection{4. An exact rank-five nonsquare
example}\label{SF-E839:an-exact-rank-five-nonsquare-example}

Take the two standard bivectors in the \(z\)-coordinates and put
\(\ell=Jz\), where

\[
 J=\begin{pmatrix}
 0&1&0&1\\
 1&1&3&2\\
 1&0&1&0\\
 1&0&2&0
 \end{pmatrix},\qquad \det J=-1.
 \tag{E839.10}
\]

Use the actual quadrics

\[
 F_1=\ell_1^2+\ell_2\ell_3,\qquad
 F_2=\ell_2^2,\qquad F_3=\ell_3^2,\qquad F_4=\ell_4^2.
 \tag{E839.11}
\]

The quotient is free with basis \(1,\ell_1\) over the length-eight
square algebra in \(\ell_2,\ell_3,\ell_4\); hence it has length sixteen
and is an Artinian quadratic complete intersection. The transformed
bivector coefficients, in pair order \(12,13,14,23,24,34\), are

\[
 p=(-1,-1,-1,-1,-1,0),\qquad
 q=(-3,-1,-2,-2,-4,0).
 \tag{E839.12}
\]

Both are decomposable and their mixed wedge equals \(\det J=-1\), as
follows directly by applying \(J\) to the two standard decomposable
forms.

Use the six pair monomials as the basis of \(S_2/Q\). The complete
unrestricted-coefficient calculation (E838.21) applies to the same four
relation shapes. With \(P_{ij}=ap_{ij}+bq_{ij}\), its six-by-six bracket
matrix is

\[
 \begin{pmatrix}
 4P_{12}&0&0&0&0&0\\
 0&4P_{13}&0&0&0&0\\
 0&0&4P_{14}&0&0&0\\
 0&0&0&4P_{23}&0&0\\
 0&0&2P_{34}&0&4P_{24}&0\\
 0&0&2P_{24}&0&0&4P_{34}
 \end{pmatrix}.
 \tag{E839.13}
\]

Here \(P_{34}=0\) identically, so the last column vanishes and every
pencil value has rank at most five. Deleting its last row and last
column gives the exact minor

\[
 -1024(a+3b)(a+b)(a+2b)^2(a+4b).
 \tag{E839.14}
\]

It is nonzero at both pure endpoints. Thus both maps \(T_{1,0}\),
\(T_{0,1}\) have rank exactly five, and the generic pencil rank is five.
In particular the entire determinant pencil of E838 is identically zero,
while the new criterion excludes this example and every higher-order
perturbation having these quadratic jets.

The example lies outside E837. Its relation space has only three
projective rank-one quadrics: a nonzero coefficient at \(F_1\) produces
an isolated \(\ell_1^2\) term and a nonzero \(\ell_2\ell_3\) term, which
cannot be a square of a linear form because all three \(\ell_1\ell_j\)
coefficients vanish. With that coefficient zero, only the three
coordinate-square lines are rank one. A formal square relation space has
four independent rank-one quadrics. Formal parameter changes and
abstract local-algebra isomorphisms preserve this cotangent-space
relation invariant.

It also lies outside E836's commuting-triple criterion for every minimal
regular generator presentation, for either original bivector. Any such
presentation changes the leading quadrics by an invertible constant
matrix. The kernel of either six-dimensional bracket map has dimension
one. Three independent leading generators have a three-dimensional
exterior-square subspace; it cannot be contained in a one-dimensional
kernel. Their three brackets therefore cannot all vanish modulo the
ideal. The same argument applies to every higher-order perturbation with
these jets.

\subsection{5. Actual geometric scope and source
boundary}\label{SF-E839:actual-geometric-scope-and-source-boundary}

The global passage remains the precise reviewed E819 one: invariant
vanishing for the specified coherent finite action implies ordinary
vanishing after forgetting the action, and completion supplies the two
local deformations. An algebraic rational \(B\) shear preserves their
bivector components. No chosen local generator basis needs an
equivariant structure.

There are actual balanced coherent objects with the new excluded local
algebra. Realize its length-sixteen finite quotient at a point of the
abelian fourfold, take the rank-thirty-two bundle \(H\) with full
character \(32v_{1/2}\), choose a quotient
\(H\twoheadrightarrow\mathcal O_Z^m\) for \(1\le m\le32\), and let \(K\)
be its kernel. Any extension \(0\to\mathcal O_Z^m\to E\to K\to0\) has
precisely this torsion, locally free hull \(H\), and identical hull
quotient. Its full character is \(32v_{1/2}\), and it is ordinarily
gluable as a coherent sheaf. The Koszul calculation gives
extension-space dimension \(96m^2\). The theorem excludes weak
semiregularity; it asserts no action, simplicity, or successful seed.

No new outside theorem is used. E838's complete associated-graded and
Jacobian argument, with the precisely located published Eisenbud--Ulrich
Theorem 7.1 and Appendix proofs, is retained as a checked input. The new
all-size matrix proof, actual rank-five witness, residual-minor count
and strict-scope statements are proved above. The finite supplement
checks exact matrices and minors only. Rank-at-most-four residual jets,
higher initial degrees, other Artinian defect modules and unrelated
perfect complexes remain open in this result.
